% ============================================================
% ［架空の記入例・提出不可］研究設定と本研究の位置づけは説明用である．
% 引用する3件の文献は実在する．
% 通常のchapters/02-background.texには自動で読み込まれない．
% ============================================================
\chapter{背景・関連研究}

\begin{center}
  \UsamiFictionalExampleNotice
\end{center}

\section{2次元人物姿勢推定}

2次元人物姿勢推定では，画像上の人物ごとに複数の関節座標を求め，所定の接続関係に従って骨格を構成する．代表的な構成には，人物領域を検出した後で個人ごとに関節を推定するtop-down方式と，画像全体から関節候補を検出して人物単位に対応づけるbottom-up方式がある．Part Affinity Fields\cite{cao2017partaffinity}は，関節候補に加えて関節間の結びつきを表す場を推定し，複数人物の骨格を構成するbottom-up方式である．

HRNet\cite{sun2019hrnet}は，解像度の異なる表現を並列に保持しながら情報交換することで，位置情報を必要とする姿勢推定に高解像度表現を利用する．ViTPose\cite{xu2022vitpose}は，Vision Transformerを基盤とした簡潔な姿勢推定の構成を示している．本研究は，これらの推定器の内部構造を変更するものではない．推定器から人物ごとの関節座標と関節信頼度を取得できることだけを前提とし，出力後の骨格を処理対象とする．

\section{低照度環境における誤推定}

低照度画像では，センサ雑音とコントラスト低下により関節周辺の局所特徴が不明瞭になる．さらに，露光時間の増加による動きぼけや，暗部での人物領域検出の欠落が重なると，関節位置だけでなく人物と関節の対応も不安定になる．その結果，身体の一部だけが別人物へ接続された骨格，胴体に対して手足が過度に長い骨格，複数関節が同一点付近へ縮退した骨格が生じ得る．

関節信頼度の平均に閾値を設ける方法は簡便であるが，各関節の局所的な確からしさしか利用しない．一方，人体骨格は関節間の接続関係をもち，同一人物内の正規化辺長は一定の範囲に収まると考えられる．そこで本例では，関節信頼度を局所的指標，骨格辺長の基準からの逸脱を構造的指標として分け，両者を統合して骨格単位の利用可否を判定する．

\section{本研究の位置づけ}

Part Affinity Fields\cite{cao2017partaffinity}が関節候補から人物骨格を組み立てる際の対応を扱うのに対し，本研究は既に構成された骨格を採用するか除外するかを判定する．また，HRNet\cite{sun2019hrnet}やViTPose\cite{xu2022vitpose}のような姿勢推定モデルを低照度データで再学習するのではなく，モデルの出力だけで完結する後処理を対象とする．このため導入は容易である一方，除外された骨格を復元できず，姿勢推定器自体の位置誤差も改善しない．本例では，この適用範囲を越えて性能を主張せず，採用骨格の精度と欠損量を同時に評価する．

本章の低照度条件，想定する誤推定および研究上の差分は，論文構成を示すための架空設定である．
